\documentclass[10pt,a4paper]{amsart}
\usepackage[margin=19mm]{geometry}
\usepackage{amsmath,amssymb,mathtools}
\usepackage[scr=rsfs,cal=euler]{mathalfa}
\usepackage{xfrac}
\usepackage[unicode,hidelinks,bookmarks=false]{hyperref}
\newcommand{\ie}{i.e.\ }
\newcommand{\cf}{cf.\ }
\newcommand{\resp}{respectively\ }
\newcommand{\Subsection}{\subsection}
\providecommand{\nobreakdash}{\nobreak\hbox{-}}
\setlength{\emergencystretch}{2em}
\multlinegap0pt
\usepackage{amssymb}
\usepackage{mathtools}
\usepackage[shortlabels]{enumitem}
\usepackage{xcolor}
\usepackage{tikz}
\newtheorem{theorem}{Theorem}
\newtheorem{lemma}{Lemma}
\newcommand{\m}{\mathfrak m}
\newcommand{\one}{\mathbf 1}
\newcommand{\rA}{\mathfrak r}
\DeclareMathOperator{\rad}{rad}
\title{Obstructions to Jacobi-Finiteness of Quivers with Potentials}
\author{Wen Chang, QuanYu Tang}
\date{Source: arXiv:2609.11200v1 (CC BY 4.0)}
\begin{document}
\maketitle

\noindent\emph{Source and licence.} Obstructions to Jacobi-Finiteness of Quivers with Potentials. Version arXiv:2609.11200v1 (CC BY 4.0), distributed by arXiv under the Creative Commons Attribution 4.0 International licence, http://creativecommons.org/licenses/by/4.0/, as recorded in the arXiv licence metadata for this exact version and shown on the abstract page https://arxiv.org/abs/2609.11200v1. Full licence notice: \texttt{licenses/arxiv-2609.11200-CC-BY-4.0.txt}.\par\smallskip

\noindent\textit{Editorial scope.} Counted unit: the article's theorem theorem:completed-matrix-gs (source0.tex lines 240--257). The article's complete original proof of it is delivered with it (source0.tex lines 591--718, one \texttt{proof} environment of 128 lines, which the article places away from the statement). No mathematics is written, repaired or re-derived: the statement and the proof are the article's own lines, in the article's own notation, and no retained line is substituted or re-flowed.  Retained uncounted as the source-local context the proof needs: lines 397--397; lines 412--420; lines 497--500; lines 536--543.  Nothing was omitted from any retained span, so the package carries no omission record.  No retained line contains a citation, so the wrapper carries no bibliography and no bibliography entry.  No retained line contains a figure environment, an image-inclusion command or drawing code, so no image is delivered and the package declares no asset dependency.  Editorial formatting only, in addition: an AMS \texttt{amsart} preamble that restates the article's own declarations and macros used by the retained lines, namely the environments \texttt{theorem}, \texttt{lemma} on separate counters, together with the standard packages those lines need.  The retained environments print in this document as Theorem 1, Lemma 1 in order of appearance, whereas the article prints the same items with the numbering of its own full document.  The complete native source of the article as distributed in the arXiv e-print for this exact version is bundled unchanged as \texttt{sources/209993/source0.tex}.
	\subsection{A filtered spanning estimate}\label{subsec:filtered-spanning-estimate}

	\begin{lemma}
		\label{lemma:filtered-spanning}
		Suppose that a finite family $(z_\lambda)_{\lambda\in\Lambda}$ spans $X$,
		and that $z_\lambda\in F^{d_\lambda}X$ for integers $d_\lambda\geq0$.
		Then
		\[
		h_\rho(X)\leq\sum_{\lambda\in\Lambda}\rho^{d_\lambda}.
		\]
	\end{lemma}

	\begin{lemma}
		\label{lemma:radical}
		If $A$ is finite-dimensional, then $\rA=\rad A$ and $\rA$ is nilpotent.
	\end{lemma}

	\begin{lemma}
		\label{lemma:kernel-generators}
		For every $\lambda$, one has
		$G_\lambda\in e_{i_\lambda}Le_{j_\lambda}$, and
		\begin{equation*}
			L=\sum_{\lambda=1}^mG_\lambda A.
		\end{equation*}
	\end{lemma}

	\begin{theorem}
		\label{theorem:completed-matrix-gs}
Let $Q$ be a finite quiver with $n$ vertices, and let $M$ be its
adjacency matrix. Let
$A=\widehat{KQ}/\mathcal I$ be defined by finitely many
endpoint-homogeneous relations $r_1,\ldots,r_m$ in
$\mathfrak m^2$.
For $0<\rho<1$, let $R(\rho)$ be
		the matrix recording the $\mathfrak m$-adic orders and the two endpoint
		idempotents of these relations.  If $A$ is finite-dimensional, then there
		exists $H(\rho)\in\mathbb R_{\geq1}^{n}$ such that
		\begin{equation}
			\label{eq:completed-matrix-gs}
			\one\leq
			\bigl({\bf I}_{n}-\rho M^{\top}+R(\rho)\bigr)H(\rho),
		\end{equation}
		where $\bf 1$ is a column vector with entries one, and $\bf I_n$ is the $n\times n$ identity matrix. 
	\end{theorem}

	\begin{proof}[Proof of Theorem~\ref{theorem:completed-matrix-gs}]
		We use the shifted filtrations
		\[
		F^0E=F^1E=E,
		\qquad
		F^qE=V\otimes_S\rA^{q-1}\quad(q\geq2),
		\]
		and
		\[
		F^0\rA=F^1\rA=\rA,
		\qquad
		F^q\rA=\rA^q\quad(q\geq2).
		\]
		Give $L$ the induced filtration $F^qL=L\cap F^qE$.  Then the multiplication map
		$\mu:E\to\rA$ is strict.  In fact for $q\geq2$,
		\[
		\mu(F^qE)=\pi(V)\rA^{q-1}
		=\pi(V\m^{q-1})=\pi(\m^q)=\rA^q,
		\]
		and strictness for $q=0,1$ is the surjectivity of $\mu$.  Thus the
		associated graded sequence of
		$0\to L\to E\overset{\mu}{\to}\rA\to0$ is exact.
		Furthermore, for each vertex $i$, the induced graded sequence 
$0\to e_iL\to e_iE\overset{\mu}{\to}e_i\rA\to0$ is also exact.

We define
\[
H_i(\rho)=\sum_{q\geq0}
\dim_K(e_i\rA^q/e_i\rA^{q+1})\rho^q,
\qquad \rA^0=A.
\]
The sum is finite by Lemma~\ref{lemma:radical}, and $H_i(\rho)\geq1$
because $A/\rA\cong S$. For $e_i\rA$, the shifted filtration
$F^0\rA=F^1\rA=\rA$ and $F^q\rA=\rA^q$ for $q\geq2$ gives a vanishing
$q=0$ layer, while every subsequent layer matches the corresponding
term of~$H_i(\rho)$.  Since the omitted $q=0$ term of~$H_i(\rho)$ is
$\dim_K(e_iA/e_i\rA)=1$, we obtain
\[h_\rho(e_i\rA)=H_i(\rho)-1,
\]
where $h_\rho$ is defined as in
Subsection~\ref{subsec:filtered-spanning-estimate}.

Since $a\in e_{t(a)}Ve_{s(a)}$, the space $e_iVe_j$ is spanned by the
arrows from $j$ to~$i$, so $\dim_K(e_iVe_j)=m_{ji}$.
Using $1=\sum_{j\in Q_0}e_j$, the $i$-th vertex component of
$E=V\otimes_SA$ decomposes as
\[
 e_iE=\bigoplus_{j\in Q_0}(e_iVe_j)\otimes_K(e_jA).
\]
We now read off $h_\rho(e_iE)$ layer by layer, using the shifted
filtration $F^0E=F^1E=E$ and $F^qE=V\otimes_S\rA^{q-1}$ for $q\geq2$.

The layer $q=0$.
Since $F^0E=F^1E$, the quotient $F^0(e_iE)/F^1(e_iE)$ vanishes.

The layer $q=1$.
We have
\[
 F^1(e_iE)/F^2(e_iE)
 \cong\bigoplus_{j\in Q_0}(e_iVe_j)\otimes_K(e_jA/e_j\rA).
\]
Since 
$e_jA/e_j\rA\cong Ke_j$ is one-dimensional, we have
\[
 \dim_K\bigl(F^1(e_iE)/F^2(e_iE)\bigr)=\sum_{j\in Q_0}m_{ji},
\]
where $\dim_K(e_iVe_j)=m_{ji}$ by our convention.

The layers $q\geq2$.
Similarly,
\[
 F^q(e_iE)/F^{q+1}(e_iE)
 \cong\bigoplus_{j\in Q_0}(e_iVe_j)\otimes_K
       (e_j\rA^{q-1}/e_j\rA^q),
\]
so
\[
 \dim_K\bigl(F^q(e_iE)/F^{q+1}(e_iE)\bigr)
 =\sum_{j\in Q_0}m_{ji}\,
   \dim_K(e_j\rA^{q-1}/e_j\rA^q).
\]

Assembling these dimensions into $h_\rho$ and factoring out
$\rho\sum_jm_{ji}$ yields
\[
 h_\rho(e_iE)
 =\rho\sum_{j\in Q_0}m_{ji}
  \Bigl[\dim_K(e_jA/e_j\rA)
   +\sum_{q\geq1}\rho^{q}\,
    \dim_K(e_j\rA^{q}/e_j\rA^{q+1})
  \Bigr]
 =\rho\sum_{j\in Q_0}m_{ji}\,H_j(\rho).
\]

Then the exactness gives
		\begin{equation}
			\label{eq:exact-weight}
			h_\rho(e_iL)
			=\rho\sum_{j\in Q_0}m_{ji}H_j(\rho)-H_i(\rho)+1.
		\end{equation}
		
		For each $j$, choose a basis $\mathcal B_j$ of $e_jA$ by lifting bases of
		the successive quotients $e_j\rA^q/e_j\rA^{q+1}$.  If
		$b\in\mathcal B_j$ is chosen at level $q$, write $\nu(b)=q$.  Then
		\[
		H_j(\rho)=\sum_{b\in\mathcal B_j}\rho^{\nu(b)}.
		\]
		
		The definition of $d_\lambda$ gives
		$G_\lambda\in F^{d_\lambda}E$.  By
		Lemma~\ref{lemma:kernel-generators}, the vectors $G_\lambda b$ with
		$i_\lambda=i$ and $b\in\mathcal B_{j_\lambda}$ span $e_iL$.
		Such a vector belongs to $F^{d_\lambda+\nu(b)}L$.
		Lemma~\ref{lemma:filtered-spanning} implies
		\begin{equation}
			\label{eq:relation-weight}
			h_\rho(e_iL)
			\leq\sum_{j\in Q_0}R(\rho)_{ij}H_j(\rho).
		\end{equation}
		
		Combining \eqref{eq:exact-weight} and \eqref{eq:relation-weight} gives
		\[
		1\leq H_i(\rho)-\rho\sum_{j\in Q_0}m_{ji}H_j(\rho)
		+\sum_{j\in Q_0}R(\rho)_{ij}H_j(\rho)
		\qquad(i\in Q_0).
		\]
		This is precisely \eqref{eq:completed-matrix-gs}.
	\end{proof}

\end{document}
